\clearpage\thispagestyle{empty}
\pdfbookmark[0]{Sources and verse translations for the Mollweide review}{mollweide1814-vi-sources}
\begin{center}\Large\bfseries Sources and verse translations for the Mollweide review\end{center}

This AI translation follows the review by Carl Friedrich Gauss in \emph{Werke}, vol.~VI (Göttingen, 1874), pp.~567--569. The first printing appeared in \emph{Göttingische gelehrte Anzeigen}, no.~71, 2 May 1814, pp.~705--710. The collected review in Volume VI ends after the first essay and refers to Volume IV (1873), pp.~361--362, for the other essays and the letters. The forms Fomahand and Fohamond follow the collected text.

Karl Brandan Mollweide’s \emph{Commentationes mathematico-philologicae tres} was published in Leipzig by Carl Cnobloch in 1813. The copy linked below is held by ETH-Bibliothek Zürich (Rar~21664). Its first essay quotes the passage from \emph{Georgics} IV on p.~4 and discusses alternative readings, including Samuel Horsley’s, on pp.~28--30.

\noindent\textbf{The title cited in the review.} Three mathematical and philological essays explaining two difficult passages, one in Virgil and the other in Plato, and examining two rules of measurement given by Columella. The author of these essays, Karl Brandan Mollweide, professor of astronomy at the University of Leipzig, has added a letter to the distinguished J.~G. Schneider on the geometrical excerpts of Epaphroditus and Vitruvius Rufus.

\noindent\textbf{Virgil, \emph{Georgics} IV.231--235.}
\begin{quote}\small Twice they gather a heavy yield. There are two harvest times,\\
as soon as Taygete the Pleiad shows her lovely face to the earth\\
and with her foot thrusts back the waters of Ocean she spurns,\\
or when, fleeing the star of the watery Fish,\\
she descends, sadder, from the sky into the winter waves.\end{quote}

\noindent\textbf{Virgil, \emph{Georgics} I.215--218.}
\begin{quote}\small Spring is the time to sow beans. Then the loose furrows\\
receive you too, lucerne, and the yearly care for millet begins,\\
when the shining Bull opens the year with his gilded horns\\
and the Dog sets, \emph{yielding} to the opposing star.\end{quote}

First printing (\href{https://images.p.sub.uni-goettingen.de/iiif/image/gdz:PPN319721507_1814:00000709/full/1650,/0/default.jpg}{705}, \href{https://images.p.sub.uni-goettingen.de/iiif/image/gdz:PPN319721507_1814:00000710/full/1650,/0/default.jpg}{706}, \href{https://images.p.sub.uni-goettingen.de/iiif/image/gdz:PPN319721507_1814:00000711/full/1650,/0/default.jpg}{707}, \href{https://images.p.sub.uni-goettingen.de/iiif/image/gdz:PPN319721507_1814:00000712/full/1650,/0/default.jpg}{708}).

\href{https://www.e-rara.ch/zut/content/titleinfo/15849974}{Book catalogue}, \href{https://www.e-rara.ch/i3f/v20/15849980/full/1400,/0/default.jpg}{Title page}, \href{https://www.e-rara.ch/i3f/v20/15849989/full/1400,/0/default.jpg}{Virgil in Mollweide, p.~4}. Mollweide’s conclusions and alternative readings, pp.~\href{https://www.e-rara.ch/i3f/v20/15850013/full/1400,/0/default.jpg}{28}, \href{https://www.e-rara.ch/i3f/v20/15850014/full/1400,/0/default.jpg}{29}, \href{https://www.e-rara.ch/i3f/v20/15850015/full/1400,/0/default.jpg}{30}.

\clearpage
